% Informationen ------------------------------------------------------------
% 	Definition von globalen Parametern, die im gesamten Dokument verwendet
% 	werden können (z.B auf dem Deckblatt etc.).
% --------------------------------------------------------------------------
\newcommand{\titel}{Autonomes Landen von Drohnen in
Agrarlandschaften basierend auf
Deep-Reinforcement-Learning}
\newcommand{\art}{Bachelorarbeit}
\newcommand{\ort}{Leipzig}
\newcommand{\hochschule}{Universität Leipzig}
\newcommand{\fachgebiet}{Professur Informationsmanagement}
\newcommand{\fakultaet}{Wirtschaftswissenschaftliche Fakultät}
\newcommand{\institut}{Institut für Wirtschaftsinformatik}
\newcommand{\erstbetreuer}{Prof. Dr. Bogdan Franczyk}
\newcommand{\zweitbetreuer}{Nico Heider}

\newcommand{\autor}{Fritz Körner}
\newcommand{\matrikelnr}{3748012}
\newcommand{\email}{fk67rahe@studserv.uni-leipzig.de}
\newcommand{\telefonnummer}{+4915221457367}
\newcommand{\strasse}{Klarastraße 17}
\newcommand{\ortplz}{04229 Leipzig}

\newcommand{\eingereicht}{15. Juni 2026}

% Grad/Studiengang – anpassen, falls nötig
\providecommand{\grad}{Bachelor of Science}
\providecommand{\studiengang}{Informatik}


% Eigene Befehle
\newcommand{\todo}[1]{\textbf{\textsc{\textcolor{red}{(TODO: #1)}}}}

% Autorennamen in small caps
\newcommand{\AutorZ}[1]{\textsc{#1}}
\newcommand{\Autor}[1]{\AutorZ{\citeauthor{#1}}}

% Befehle zur semantischen Auszeichnung von Text
\newcommand{\Prozess}[1]{\textit{#1}}
\newcommand{\Webservice}[1]{\textit{#1}}
\newcommand{\Eingabe}[1]{\texttt{#1}}
\newcommand{\Code}[1]{\texttt{#1}}
\newcommand{\Datei}[1]{\texttt{#1}}
\newcommand{\Datentyp}[1]{\textsf{#1}}
\newcommand{\XMLElement}[1]{\textsf{#1}}

% Abkürzungen
\newcommand{\vgl}{Vgl.\ }
\newcommand{\ua}{\mbox{u.\,a.\ }}
\newcommand{\zB}{\mbox{z.\,B.\ }}
\newcommand{\bs}{$\backslash$}

% Einfache Anführungszeichen in texttt
\newcommand{\sq}{\textquotesingle}


